\documentclass[11pt]{article}
\usepackage[margin=1in]{geometry}
\usepackage{amsmath,amssymb,amsthm}
\usepackage[colorlinks=true,urlcolor=blue]{hyperref}
\newtheorem{theorem}{Theorem}
\newtheorem{proposition}[theorem]{Proposition}
\newtheorem{remark}[theorem]{Remark}
\newcommand{\one}{\boldsymbol 1}
\newcommand{\tr}{\operatorname{tr}}
\title{One doubled edge can maximally change the unique optimal bisection\\
while its entire relative fractional response vanishes}
\author{Partitioning-limits research session: elementary obstruction}
\date{November 15, 2026}
\begin{document}
\maketitle

\paragraph{Scope.}
This is a fully explicit obstruction for weighted combinatorial Laplacians.
It is not a claim of historical novelty or a field-scale breakthrough.
The host is sparse, planar, bipartite, and has maximum combinatorial degree
four. Its edge weights have growing dynamic range. Both endpoint Fiedler
eigenvalues are simple; both eigenvectors are perfectly flat and have no
zero entries; their sign cuts are the unique globally optimal cardinality
bisections. Consequently the phenomenon does not arise from an algorithmic
tie or an entrywise rounding ambiguity.

\begin{theorem}\label{thm:main}
For every integer $m\ge2$ there are connected simple weighted graphs
$G_m,G'_m$ on $n=4m$ vertices such that:
\begin{enumerate}
\item $G_m$ is planar and bipartite, has $2n-4$ edges and maximum
combinatorial degree four. The graph $G'_m$ is obtained by doubling the
conductance of one existing edge.
\item Writing $L,L'$ for their combinatorial Laplacians,
\[
             L\preceq L'\preceq(1+1/m)L.
\]
More precisely, if $L'=L+a rr^T$ is the edge update, then
$\tau=a r^TL^\dagger r<1/m$ and
$L'\preceq(1+\tau)L$.
\item The Fiedler eigenvalues of both graphs are simple. Their unit
Fiedler vectors $x,y$ satisfy
\[
 |x_i|=|y_i|=n^{-1/2},\qquad x^Ty=0.
\]
Thus $\|xx^T-yy^T\|_2=1$.
\item The sign partitions $S_x,S_y$ are balanced and are the unique
minimum-weight bisections, up to complementation, of their respective
graphs. They have the maximum possible disagreement modulo complementation:
\[
 \min\{|S_x\triangle S_y|,|S_x\triangle S_y^c|\}=n/2.
\]
\item For every $0<\alpha<1$, the relative fractional correction
\[
 C_\alpha=(L^\dagger)^{\alpha/2}
                  ((L')^\alpha-L^\alpha)(L^\dagger)^{\alpha/2}
\]
satisfies
\[
 0\preceq C_\alpha\preceq ((1+1/m)^\alpha-1)I,
 \qquad
 \tr\log(I+C_\alpha)=\alpha\log(1+\tau)
                             <\alpha\log(1+1/m).
\]
All inverse powers and determinants here act on $\one^\perp$.
\end{enumerate}
\end{theorem}

\begin{proof}
Use vertices $(s,t,j)\in\{-1,1\}^2\times\{1,\ldots,m\}$.
Set
\[
        a=\frac{4m}{4m+1},\qquad M=2m^3.
\]
The edge that flips $s$ has weight $a$, the edge that flips $t$ has
weight one, and each edge $(s,t,j)(s,t,j+1)$ has weight $M$.
This is $C_4\mathbin\square P_m$ with anisotropic weights. Draw its
$m$ four-cycles as concentric squares to see planarity. Its bipartition
is the parity of the two sign bits and $j$. There are $4m$ cycle edges
and $4(m-1)$ path edges. End layers have degree three and interior
layers degree four.

Let $H=M L(P_m)$. Its first positive eigenvalue obeys
\[
 \Lambda=4M\sin^2\!\frac{\pi}{2m}
       \ge\frac{4M}{m^2}=8m,
\]
using $\sin u\ge2u/\pi$ on $[0,\pi/2]$.
The product spectrum is formed by adding eigenvalues of $H$ to
\[
                 0,\quad2a,\quad2,\quad2a+2.
\]
In particular the vectors
\[
 x(s,t,j)=\frac{s}{\sqrt{4m}},\qquad
 y(s,t,j)=\frac{t}{\sqrt{4m}},\qquad
 z(s,t,j)=\frac{st}{\sqrt{4m}}
\]
are orthonormal eigenvectors with eigenvalues $2a,2,2a+2$.
The first is the simple Fiedler vector.

Fix $j_0$ and let $r=e_{(1,1,j_0)}-e_{(-1,1,j_0)}$. Double that
edge, so $L'=L+a rr^T$. Then
\[
 r^Ty=0,\qquad r^Tx=r^Tz=m^{-1/2}.
\]
Consequently $y$ remains an eigenvector of $L'$ with eigenvalue two.
We now prove it is the simple Fiedler vector of $L'$.

On $W=\one^\perp\cap y^\perp$, decompose into $\operatorname{span}\{x\}$
and its orthogonal complement $K$. Since
$1-a=a/(4m)$, the restriction of $L-2I$ has block form
\[
    \begin{pmatrix}-a/(2m)&0\\0&D\end{pmatrix},\qquad D\succ0.
\]
The vector $z\in K$ has $D$-eigenvalue $2a$, while all remaining
components of $r$ lie in the product modes with positive $H$ eigenvalue
and hence have $D$-eigenvalues at least $\Lambda-2$.
Writing $r=m^{-1/2}x+q$ on $W$, set $s_0=q^TD^{-1}q$.
The spectral decomposition gives
\[
 a s_0\le\frac1{2m}+\frac{2a}{\Lambda-2}
       \le\frac1{2m}+\frac2{8m-2}<1 \quad(m\ge2).
\]
The lower-right block $D+aqq^T$ is positive definite. The rank-one
inverse formula gives the Schur complement of this block in
$(L'-2I)|_W$ as
\[
 -\frac a{2m}+\frac{a/m}{1+a s_0}
       =\frac am\left(\frac1{1+a s_0}-\frac12\right)>0.
\]
Thus $L'-2I$ is positive definite on $W$, proving the claim.

For any cardinality bisection $S$, its sign vector $\sigma_S\in\{-1,1\}^n$
has mean zero, squared norm $n$, and
$\sigma_S^TL\sigma_S=4w_G(\partial S)$. Therefore
\[
       w_G(\partial S)\ge\frac{n\lambda_2(L)}4.
\]
The sign vector of $x$ attains equality, and simplicity shows that it is
the only one up to sign. The same argument applies to $y$ and $L'$.
Their cuts are $s=1$ and $t=1$, and each of their four intersections
has exactly $m$ vertices, giving the claimed distance. Explicitly, the
old $x$-cut has weight $2ma=2m-a/2$, its edited weight is $2m+a/2$,
and the $y$-cut has weight $2m$ in both graphs.

It remains to prove that the edge perturbation is small in relative
energy. Decompose $r=m^{-1/2}x+m^{-1/2}z+h$, where $h$ lies in the
positive $H$ modes. Then
\begin{align*}
 \tau=a r^TL^\dagger r
 &\le\frac1{2m}+\frac{a}{2(a+1)m}+\frac{2a}{\Lambda}\\
 &<\frac1{2m}+\frac1{4m}+\frac1{4m}=\frac1m.
\end{align*}
Cauchy--Schwarz in the $L$ energy shows
$a rr^T\preceq\tau L$, proving the relative order statement.
For $0<\alpha<1$, operator monotonicity of the fractional power gives
\[
    L^\alpha\preceq(L')^\alpha\preceq(1+\tau)^\alpha L^\alpha.
\]
Congruence by $(L^\dagger)^{\alpha/2}$ gives the operator bound.
Finally, on $\one^\perp$ the determinant lemma gives
$\det L'/\det L=1+\tau$, and hence
\[
 \det(I+C_\alpha)
       =\frac{\det((L')^\alpha)}{\det(L^\alpha)}=(1+\tau)^\alpha.
\]
Taking logarithms proves the last identity.
\end{proof}

\begin{proposition}[Exact gap-free worst-case law]
For every $\varepsilon>0$, within the class of connected simple
weighted planar bipartite graphs of maximum degree four, even when
both endpoint Fiedler vectors are flat and their sign cuts are the
unique optimal bisections, the supremum of the normalized bisection
distance under one edge doubling with
$L\preceq L'\preceq(1+\varepsilon)L$ equals $1/2$.
The corresponding supremum of the Fiedler projector distance equals
one. Both suprema are attained by a finite graph.
\end{proposition}
\begin{proof}
The universal upper bounds follow from taking the smaller of a cut
distance and its complement, and from the distance between orthogonal
projectors of rank one. Choose $m\ge\max\{2,1/\varepsilon\}$ in
Theorem~\ref{thm:main} to attain both bounds.
\end{proof}

\begin{proposition}[A necessary parameter escape for vanishing leverage]
If a host has maximum combinatorial degree $d$ and every conductance
lies in $[c,C]$, then doubling an existing edge of conductance $a$ has
relative perturbation norm
\[
                 \tau=a r^TL^\dagger r
                          \ge\frac{2c}{(d+1)C}.
\]
In particular, vanishing relative energy change under a single edge
doubling requires unbounded degree or unbounded conductance ratio.
\end{proposition}
\begin{proof}
Cauchy--Schwarz in the spectral decomposition on $\one^\perp$ gives
\[
 (r^TL^\dagger r)(r^TLr)\ge\|r\|^4=4.
\]
For an existing edge $uv$ of conductance $a$,
$r^TLr=d_w(u)+d_w(v)+2a\le2(d+1)C$, while $a\ge c$.
\end{proof}

\begin{proposition}[Only two exact balanced-cut envelopes]
Let a graph have one edge conductance increased by $t\ge0$, and impose
any fixed family $\mathcal F$ of admissible cuts, such as cardinality
bisections. Among cuts that separate the edge endpoints, let $A$ be the
minimum original cut weight and retain all cuts attaining $A$; among
cuts that do not separate them, define $B$ similarly. The optimum
throughout the whole trajectory is
\[
                          \min\{A+t,B\}.
\]
Its complete optimizer set is one of these two fixed cut families,
or their union at the unique possible crossing. Empty classes have
minimum $+\infty$.
\end{proposition}
\begin{proof}
Each cut has cost $c(S)+t\one_{\{e\in\partial S\}}$. Within each of
the two classes, adding the same quantity does not change its complete
minimizer set.
\end{proof}

\paragraph{A four-vertex variant with comparable weights.}
For arbitrarily small changes of an edge weight, growing dimension and
large weight ratios are unnecessary. On $C_4$ give the two $s$ edges
weight $a=1-\delta/4$, the two $t$ edges weight one, and increase one
$s$ edge by $\delta$, where $0<\delta\le1/4$. The original eigenvalues
are $0,2a,2,2a+2$. The updated $x,z$ block has eigenvalues
\[
       3+\delta/2\ \pm\sqrt{1+\delta^2}.
\]
Its smaller eigenvalue exceeds two for $0<\delta<4/3$.
Thus the flat Fiedler vector changes exactly from $x$ to $y$, and their
unique optimal bisections have distance $n/2$. All weights lie in
$[15/16,5/4]$, while
$L\preceq L'\preceq(1+\delta/a)L$.

\paragraph{Prior-art boundary.}
The necessity of a spectral separation in eigenvector perturbation
theory is classical. Andreotti--Edelmann--Guglielmi--Lubich introduce a
graph-structured distance to vanishing eigengap
\cite{AEGL}; Peng--Yoshida study average sensitivity under random edge
deletions \cite{PY}; Noschese--Reichel explicitly analyze targeted
single-edge perturbations of Fiedler eigenpairs \cite{NR}.
The elementary construction here combines an exact edge doubling,
vanishing relative perturbation, bounded combinatorial degree,
planarity, flat simple eigenvectors, and unique optimal bisections.
No inspected source establishes priority for that conjunction, and
this note makes no claim of firstness. The theorem concerns a selected
worst-case edit, not average sensitivity. It does not establish a
corresponding obstruction for normalized Laplacians, uniformly
comparable weights under edge doubling, or simple unweighted
bounded-degree hosts. The edge weights $2m^3$ are essential to the
stated sparse construction's proof of vanishing leverage.

\begin{thebibliography}{9}
\bibitem{AEGL} E.~Andreotti, D.~Edelmann, N.~Guglielmi, C.~Lubich,
Measuring the stability of spectral clustering,
\emph{Linear Algebra Appl.} 610 (2021), 673--697.
\url{https://arxiv.org/abs/1903.05193}.
\bibitem{PY} P.~Peng, Y.~Yoshida,
Average Sensitivity of Spectral Clustering, KDD 2020.
\url{https://arxiv.org/abs/2006.04094}.
\bibitem{NR} S.~Noschese, L.~Reichel,
Sensitivity of Perron and Fiedler eigenpairs to structural perturbations
of a network, 2025.
\url{https://arxiv.org/abs/2509.16024}.
\end{thebibliography}
\end{document}
